\scanpage{8}
\mspage{7}

Cette structure ne change pas si
\deleted{\unsure{de}}on change l'ordre \unsure{par} ordre opposé,
i.e. elle ne dépend que du
biordre.

\deleted{Partie modérée} \deleted{\unsure{Segments|Segment}}

% slanted
\margin{refaire, sans se borner à une ligne : \unsure{sous tronçon|son tronçon} fermé d'un tronçon. C'est un tronçon \unsure{p.|pour} le biordre induit. D'autre part, le bord d'un tronçon est défini \unsure{comme|Ens} des}

\unsure{Tronçons|Tronçon} de la ligne : parties
de la forme

\[
\left\{\begin{aligned}
 &[s,t] &&s\leq t\\
 &\mathopen{]}-\infty,s],\quad\text{ou}\quad [s,+\infty\mathclose{[}\\
 &\unsure{\mathcal L|L}=\mathopen{]}-\infty,+\infty\mathclose{[}
\end{aligned}\right.
\]

Bord du tronçon : $\{s,t\}$ dans le
premier cas, $\{s\}$ dans le second, $\varnothing$ \inserted{troisième}

Deux tronçons sont dits en
position relative modérée, si leurs
bords \unsure{:|} \unsure{le|les} sont.

\uline{Partie} \uline{modérée} :
réunion \inserted{finie} de tronçons \inserted{en position} mutuelle
modérée, et disjoints. \deleted{\unsure{Donc}} Par
\unsure{définition|definition}, son bord est la réunion
(disjointe) des bords.

% paper watermark; clipped at the two sheet edges
\rotatebox{90}{\textsc{\unsure{r}ong Extra Stro}}
